\documentclass[a4paper,12pt]{article}

% -------------------------------------------------
% Кодировка, язык, шрифт
% -------------------------------------------------
\usepackage[T2A]{fontenc}
\usepackage[utf8]{inputenc}
\usepackage[russian]{babel}
\usepackage{helvet}
\renewcommand{\familydefault}{\sfdefault}

% -------------------------------------------------
% Математика и типографика
% -------------------------------------------------
\usepackage{amsmath,amssymb,mathtools}
\usepackage{icomma}
\usepackage{microtype}
\usepackage{setspace}
\usepackage{parskip}

% -------------------------------------------------
% Страница
% -------------------------------------------------
\usepackage[
  a4paper,
  left=19mm,
  right=19mm,
  top=19mm,
  bottom=21mm
]{geometry}

\usepackage{enumitem}
\usepackage{array,tabularx,booktabs}
\usepackage{xcolor}
\usepackage{tikz}
\usetikzlibrary{calc}
\usepackage{fancyhdr}
\usepackage{titlesec}

% -------------------------------------------------
% Цвета
% -------------------------------------------------
\definecolor{accent}{HTML}{008080}
\definecolor{darkaccent}{HTML}{004D4D}
\definecolor{textgray}{HTML}{333333}
\definecolor{softgray}{HTML}{707070}
\definecolor{linegray}{HTML}{D6DEDE}

% -------------------------------------------------
% Общие настройки
% -------------------------------------------------
\color{textgray}
\onehalfspacing
\setlength{\parindent}{0pt}
\setlength{\headheight}{15pt}
\setlength{\emergencystretch}{2em}

\setlist[itemize]{
  leftmargin=6mm,
  itemsep=2pt,
  topsep=4pt
}

\setlist[enumerate]{
  leftmargin=8mm,
  itemsep=7pt,
  topsep=5pt
}

\titleformat{\section}
  {\Large\bfseries\color{darkaccent}}
  {}
  {0pt}
  {}

\titleformat{\subsection}
  {\large\bfseries\color{accent}}
  {}
  {0pt}
  {}

\titlespacing*{\section}{0pt}{15pt}{7pt}
\titlespacing*{\subsection}{0pt}{11pt}{5pt}

\pagestyle{fancy}
\fancyhf{}
\lhead{\small\color{softgray} Степени и степенные преобразования}
\rhead{\small\color{softgray} Анастасия Павлова}
\cfoot{\small\color{softgray}\thepage}
\renewcommand{\headrulewidth}{0.4pt}
\renewcommand{\headrule}{%
  \hbox to\headwidth{\color{linegray}\leaders\hrule height \headrulewidth\hfill}
}

\newcommand{\checkitem}{\(\square\)\;}
\newcommand{\important}[1]{\textcolor{darkaccent}{\textbf{#1}}}
\newcommand{\restr}[1]{\textcolor{softgray}{\small #1}}

\begin{document}

% =================================================
% ТИТУЛЬНЫЙ ЛИСТ
% =================================================
\begin{titlepage}
\thispagestyle{empty}

\begin{center}

\vspace*{28mm}

{\Large\color{accent}\textbf{ЧЕК-ЛИСТ}}

\vspace{9mm}

{\fontsize{26}{31}\selectfont
\bfseries
Степени и степенные\\[2mm]
преобразования выражений
}

\vspace{8mm}

{\large
Вычисления, корни, дробные и отрицательные показатели,\\
буквенные выражения
}

\vspace{17mm}

{\color{darkaccent}\rule{55mm}{0.8pt}}

\vspace{17mm}

{\large
\textbf{Подготовка к ЕГЭ}
}

\vfill

{\large
Лёвин Артём Александрович\\[2mm]
\textbf{эксклюзивно для Анастасии Павловой}
}

\vspace{8mm}

{\color{softgray}\today}

\vspace{23mm}

\end{center}

% Сдержанная титульная кривая Безье
\begin{tikzpicture}[remember picture,overlay]
  \draw[
    accent,
    line width=1.3pt
  ]
  ($(current page.south west)+(18mm,15mm)$)
  .. controls
  ($(current page.south west)+(63mm,26mm)$)
  and
  ($(current page.south east)+(-68mm,6mm)$)
  ..
  ($(current page.south east)+(-18mm,17mm)$);
\end{tikzpicture}

\end{titlepage}

% =================================================
% ОСНОВНОЙ ЧЕК-ЛИСТ
% =================================================
\section{1. Главная стратегия}

Почти любую вычислительную задачу со степенями удобно начинать с одного вопроса:

\begin{center}
\important{Можно ли привести составные числа, корни и дроби к степенному виду?}
\end{center}

После такого перехода выражение часто сводится к нескольким стандартным формулам.

\begin{itemize}
  \item \checkitem составное число представить степенью удобного основания;
  \item \checkitem корень заменить степенью с дробным показателем;
  \item \checkitem десятичную дробь сначала перевести в обыкновенную;
  \item \checkitem отрицательный показатель связать с обратным числом;
  \item \checkitem степени одного основания собрать вместе;
  \item \checkitem отдельно и аккуратно вычислить итоговый показатель.
\end{itemize}

\subsection{Быстрые преобразования}

\renewcommand{\arraystretch}{1.35}
\begin{table}[htbp]
\centering
\caption{Полезные переходы к степенному виду}
\begin{tabularx}{\linewidth}{@{}p{0.29\linewidth}X@{}}
\toprule
\textbf{Исходный объект} & \textbf{Полезная запись} \\
\midrule
\(49\) & \(49=7^2\) \\
\(8\) & \(8=2^3\) \\
\(18\) & \(18=2\cdot 3^2\) \\
\(0{,}5\) & \(0{,}5=\dfrac12=2^{-1}\) \\
\(0{,}2\) & \(0{,}2=\dfrac15=5^{-1}\) \\
\(\sqrt[k]{a^m}\), \(a>0\) &
\(\sqrt[k]{a^m}=a^{m/k}\) \\
\bottomrule
\end{tabularx}
\end{table}

\section{2. Формулы, которые должны работать автоматически}

\subsection{Умножение степеней одного основания}

\[
a^m\cdot a^n=a^{m+n}.
\]

\textbf{Пример:}
\[
3^{\sqrt5+10}\cdot3^{-5-\sqrt5}
=
3^{(\sqrt5+10)+(-5-\sqrt5)}
=
3^5.
\]

\important{При умножении показатели складываются.}

\subsection{Деление степеней одного основания}

При \(a\neq0\):
\[
\frac{a^m}{a^n}=a^{m-n}.
\]

Особенно внимательно следует работать со сложным вычитаемым:

\[
\frac{a^p}{a^{q-r}}
=
a^{\,p-(q-r)}
=
a^{p-q+r}.
\]

Скобки здесь защищают знаки.

\subsection{Степень степени}

\[
\left(a^m\right)^n=a^{mn}.
\]

\textbf{Пример:}
\[
\left(7^2\right)^{6\cdot3}
=
7^{2\cdot6\cdot3}.
\]

Показатели \important{перемножаются}.

\subsection{Степень произведения}

Для используемых в задачах допустимых значений:

\[
(ab)^n=a^n b^n.
\]

Например,
\[
(3x^{10})^2
=
3^2\left(x^{10}\right)^2
=
9x^{20}.
\]

Квадрат относится ко \important{всему содержимому скобок}.

\subsection{Нулевая степень}

При \(a\neq0\):
\[
a^0=1.
\]

Если после преобразований получилось
\[
x^0,
\]
следует убедиться, что исходное выражение допускает \(x\neq0\).

\subsection{Отрицательная степень}

При \(a\neq0\):
\[
a^{-n}=\frac1{a^n},
\qquad
\frac1{a^n}=a^{-n}.
\]

Например,
\[
5^{-3}
=
\frac1{5^3}
=
\frac1{125}.
\]

Знак «минус» в показателе означает переход к обратному числу.

\section{3. Корни как степени}

Для положительного основания:

\[
\sqrt[k]{a^m}=a^{m/k}.
\]

В частности,
\[
\sqrt a=a^{1/2},
\qquad
\sqrt[3]{a}=a^{1/3}.
\]

\textbf{Пример:}
\[
\frac{\sqrt[15]{5}\cdot\sqrt[10]{5}}
     {\sqrt[6]{5}}
=
5^{\frac1{15}+\frac1{10}-\frac16}.
\]

Общий знаменатель равен \(30\):

\[
\frac1{15}+\frac1{10}-\frac16
=
\frac2{30}+\frac3{30}-\frac5{30}
=
0.
\]

Следовательно,
\[
5^0=1.
\]

\subsection{ОДЗ при работе с корнями}

\begin{itemize}
  \item для корня чётной степени подкоренное выражение должно быть неотрицательным;
  \item знаменатель не может обращаться в ноль;
  \item при использовании единого степенного алгоритма с дробными показателями удобно отдельно фиксировать условие \(a>0\).
\end{itemize}

Например,
\[
\left(\frac{\sqrt a}{a^{1/3}}\right)^4
\]
требует
\[
a>0.
\]

После этого:
\[
\left(
a^{1/2-1/3}
\right)^4
=
\left(a^{1/6}\right)^4
=
a^{2/3}.
\]

\section{4. Десятичные дроби}

Десятичную дробь полезно убрать как можно раньше.

\[
0{,}5=\frac12=2^{-1},
\qquad
0{,}2=\frac15=5^{-1}.
\]

Например,
\[
(0{,}5)^{\sqrt{10}+1}
=
\left(2^{-1}\right)^{\sqrt{10}+1}
=
2^{-\sqrt{10}-1}.
\]

Так выражение сразу становится пригодным для стандартных действий со степенями.

\section{5. Скобки: защита от ошибок со знаками}

Полезная памятка занятия:

\begin{center}
\important{ВУПС: вычитаемое, умножаемое и подставляемое выражение --- в скобки.}
\end{center}

Особенно полезно применять этот принцип, когда внутри выражения присутствуют \(+\) или \(-\).

Например:
\[
a^{\,3\sqrt7-1+1-\sqrt7-(2\sqrt7-1)}.
\]

Сначала раскрываем скобки:
\[
-(2\sqrt7-1)=-2\sqrt7+1.
\]

Только затем приводим подобные слагаемые.

\section{6. Буквенные выражения}

\subsection{Приведение подобных}

Если буквенная часть полностью совпадает, коэффициенты можно сложить:

\[
7m^{30}+11m^{30}
=
18m^{30}.
\]

Тогда
\[
\frac{7m^{30}+11m^{30}}{(3m^{15})^2}
=
\frac{18m^{30}}{9m^{30}}
=
2,
\qquad m\neq0.
\]

\subsection{Сокращение степеней}

Сокращать можно множители.

Например,
\[
\frac{18x^{20}}{9x^{20}}
=
2,
\qquad x\neq0.
\]

Условие \(x\neq0\) сохраняется из исходного выражения.

\subsection{Двоеточие лучше заменить дробной чертой}

Запись
\[
18x^7\cdot x^{13}:(3x^{10})^2
\]
удобнее сразу переписать:

\[
\frac{18x^7\cdot x^{13}}{(3x^{10})^2}.
\]

Так структура числителя и знаменателя становится визуально однозначной.

\section{7. Типичные ошибки}

\begin{enumerate}[label=\textbf{\arabic*)}]
  \item
  \textbf{Применение формулы через знак сложения.}

  Из
  \[
  a^m+a^n
  \]
  нельзя получить \(a^{m+n}\).

  Формула
  \[
  a^m\cdot a^n=a^{m+n}
  \]
  связана с умножением.

  \item
  \textbf{Потеря внешней степени.}

  В выражении
  \[
  \left(\frac{\sqrt a}{a^{1/3}}\right)^4
  \]
  после преобразования содержимого скобок степень \(4\) сохраняется.

  \item
  \textbf{Степень применяется только к одному множителю.}

  Следует помнить:
  \[
  (3x^{10})^2=9x^{20}.
  \]

  \item
  \textbf{Ошибка при представлении числа степенью.}

  \[
  18=2\cdot3^2,
  \qquad
  3^3=27.
  \]

  \item
  \textbf{Потеря минуса при вычитании показателя.}

  \[
  p-(q-r)=p-q+r.
  \]

  \item
  \textbf{Преждевременная подстановка значения переменной.}

  Сначала удобнее упростить буквенное выражение, затем выполнить числовую подстановку.

  \item
  \textbf{Слишком крупные шаги решения.}

  Разделение преобразования на короткие действия облегчает самопроверку.
\end{enumerate}

\section{8. Алгоритм решения}

Перед сдачей решения проверьте каждый пункт:

\begin{enumerate}[label=\checkitem]
  \item Определены ограничения и ОДЗ.
  \item Составные числа представлены удобными степенями.
  \item Корни переведены в дробные показатели.
  \item Десятичные дроби заменены обыкновенными дробями или степенями.
  \item Сложные вычитаемые и подставляемые выражения заключены в скобки.
  \item При умножении одинаковых оснований показатели сложены.
  \item При делении одинаковых оснований показатели вычтены.
  \item При возведении степени в степень показатели перемножены.
  \item Степень произведения применена ко всем множителям в скобках.
  \item Дробные показатели приведены к общему знаменателю.
  \item Подобные слагаемые приведены.
  \item Итоговый показатель вычислен отдельно.
  \item Числовая подстановка выполнена после упрощения.
  \item Ответ проверен на соответствие исходным ограничениям.
\end{enumerate}

% =================================================
% ТРЕНИРОВОЧНЫЙ БЛОК
% =================================================
\clearpage

\section{Тренировочный блок}

\textbf{Путь Б: 04.09--06.09.}

На каждый день запланировано пять заданий.

\textbf{Режим работы с заданием 5:}
включите таймер и запишите время выполнения.
Точность решения сохраняет приоритет перед скоростью.

% -------------------------------------------------
\subsection{04.09}

\begin{enumerate}[label=\textbf{\arabic*.}]
  \item Найдите значение выражения:
  \[
  3^{\sqrt5+10}\cdot3^{-5-\sqrt5}.
  \]

  \item Найдите значение выражения:
  \[
  \frac{(49^6)^3}{(7^7)^5}.
  \]

  \item Найдите значение выражения:
  \[
  \frac{
    5^{3\sqrt7-1}\cdot5^{1-\sqrt7}
  }{
    5^{2\sqrt7-1}
  }.
  \]

  \item Найдите значение выражения:
  \[
  \frac{
    8^{1-\sqrt7}
  }{
    2^{2-3\sqrt7}
  }.
  \]

  \item Найдите значение выражения:
  \[
  \frac{
    (0{,}5)^{\sqrt{10}+1}
  }{
    2^{-\sqrt{10}}
  }.
  \]
\end{enumerate}

% -------------------------------------------------
\subsection{05.09}

\begin{enumerate}[label=\textbf{\arabic*.}]
  \item Найдите значение выражения:
  \[
  \frac{
    \sqrt[15]{5}\cdot\sqrt[10]{5}
  }{
    \sqrt[6]{5}
  }.
  \]

  \item При \(a>0\) упростите выражение:
  \[
  \left(
    \frac{\sqrt a}{a^{1/3}}
  \right)^6.
  \]

  \item При \(m\neq0\) упростите выражение:
  \[
  \frac{
    7m^{30}+11m^{30}
  }{
    (3m^{15})^2
  }.
  \]

  \item При \(x\neq0\) упростите выражение:
  \[
  \frac{
    18x^7\cdot x^{13}
  }{
    (3x^{10})^2
  }.
  \]

  \item При \(x\neq0\) найдите значение выражения:
  \[
  \frac{
    (3x)^3\cdot x^{-9}
  }{
    2x^{-10}\cdot x^4
  }.
  \]
\end{enumerate}

% -------------------------------------------------
\subsection{06.09}

\begin{enumerate}[label=\textbf{\arabic*.}]
  \item Найдите значение выражения:
  \[
  (0{,}2)^2\cdot5^{-1}.
  \]

  \item Упростите выражение:
  \[
  \frac{
    \sqrt[12]{2^5}\cdot\sqrt[8]{2^3}
  }{
    \sqrt[6]{2}
  }.
  \]

  \item При \(a\neq0\), \(b\neq0\) упростите выражение:
  \[
  \frac{
    (4a^6b^{-2})^2
  }{
    8a^5b^{-1}
  }.
  \]

  \item При \(x\neq0\) упростите выражение:
  \[
  \frac{
    6x^{12}+9x^{12}
  }{
    (3x^6)^2
  }.
  \]

  \item При \(a>0\) упростите выражение
  \[
  \left(
    \frac{\sqrt a}{a^{1/3}}
  \right)^4,
  \]
  затем найдите его значение при
  \[
  a=8.
  \]
\end{enumerate}

% =================================================
% ОТВЕТЫ
% =================================================
\clearpage

\section{Ответы к тренировочному блоку}

\renewcommand{\arraystretch}{1.35}

\begin{table}[htbp]
\centering
\caption{Ответы к заданиям 04.09--06.09}
\begin{tabularx}{\linewidth}{@{}p{0.19\linewidth}p{0.18\linewidth}X@{}}
\toprule
\textbf{Дата} & \textbf{Задание} & \textbf{Ответ} \\
\midrule

04.09 & 1 & \(243\) \\
04.09 & 2 & \(7\) \\
04.09 & 3 & \(5\) \\
04.09 & 4 & \(2\) \\
04.09 & 5 & \(\dfrac12\) \\

\midrule

05.09 & 1 & \(1\) \\
05.09 & 2 & \(a\) \\
05.09 & 3 & \(2\) \\
05.09 & 4 & \(2\) \\
05.09 & 5 & \(\dfrac{27}{2}=13{,}5\) \\

\midrule

06.09 & 1 & \(\dfrac1{125}\) \\
06.09 & 2 & \(2^{5/8}\) \\
06.09 & 3 & \(\dfrac{2a^7}{b^3}\) \\
06.09 & 4 & \(\dfrac53\) \\
06.09 & 5 & \(4\) \\

\bottomrule
\end{tabularx}
\end{table}

\vspace{8mm}

\begin{center}
\color{darkaccent}
\textbf{Финальная проверка}
\end{center}

\begin{itemize}
  \item \checkitem Все корни, дроби и составные числа рассмотрены как кандидаты на перевод в степенной вид.
  \item \checkitem При делении сложный показатель знаменателя взят в скобки.
  \item \checkitem Отрицательные показатели обработаны через обратные величины.
  \item \checkitem Внешняя степень применена ко всему содержимому скобок.
  \item \checkitem Ограничения на переменные сохранены до конца решения.
  \item \checkitem Арифметика итоговых показателей перепроверена отдельно.
\end{itemize}

\end{document}